% Generated draft: requires mathematical and source review.
\documentclass[11pt]{article}
\usepackage[margin=25mm]{geometry}
\usepackage{fontspec}
\setmainfont{Latin Modern Roman}
\usepackage{amsmath,amssymb,mathtools,unicode-math}
\setmathfont{Latin Modern Math}
\usepackage{enumitem,needspace,xurl,hyperref}
\setlength{\parindent}{0pt}
\setlength{\parskip}{5pt}
\setlength{\emergencystretch}{3em}
\newcommand{\N}{\mathbb{N}}
\newcommand{\Z}{\mathbb{Z}}
\newcommand{\Q}{\mathbb{Q}}
\newcommand{\R}{\mathbb{R}}
\newcommand{\C}{\mathbb{C}}
\newcommand{\F}{\mathbb{F}}
\newcommand{\A}{\mathbb{A}}
\newcommand{\PP}{\mathbb{P}}
\newcommand{\E}{\mathbb{E}}
\newcommand{\eps}{\varepsilon}
\newcommand{\dd}{\,\mathrm{d}}
\begin{document}
\section*{UnsolvedMath ID: 700903}
\textbf{Source problem code:} catalogue entry
\subsection{Definitions and mathematical statement}
For a nontrivial finite simple group \(G\) and an integer \(n\geq 1\), let \(G^n=\prod_{i=1}^{n}G\), with coordinatewise multiplication. A group \(H\) is called \emph{2-generator} if there exist \(x,y\in H\) such that \(H=\langle x,y\rangle\); this means that at most two generators suffice.

The problem asks, for each known finite simple group \(G\), to determine the largest positive integer \(n\) such that \(G^n\) is 2-generator. Equivalently, determine \(m\geq 1\) satisfying
\[
\bigl(\exists x,y\in G^m\bigr)\;\langle x,y\rangle=G^m,
\qquad
\bigl(\forall n>m\bigr)\bigl(\forall u,v\in G^n\bigr)\;\langle u,v\rangle\neq G^n.
\]
Here the source's word ``known'' does not specify a precise collection or a reference date. Under the usual meaning of finite simple group, cyclic groups of prime order are included; the supplied statement gives no restriction to nonabelian groups. The interpretation of ``2-generator'' above is the usual at-most-two convention; the source does not explicitly discuss whether exactly two generators must be necessary.
\subsection{Short English statement}
For every known finite simple group, determine the greatest number of copies whose direct product can be generated by two elements. The scope of ``known'' remains unspecified in the supplied source.
\subsection{Sources}
\begin{itemize}
\item Supplied local UnsolvedMath registry record: permanent ID \(700903\), exact source problem code \texttt{catalogue entry}, title \emph{Kourovka Notebook Problem 12.86}. Upstream revision: \texttt{d371ee917ac9647fbb06cd6fc529deec6f3dda82}.
\item Attributed to S. A. Syskin, proposed in 1992. The supplied background identifies \emph{The Kourovka Notebook (Unsolved Problems in Group Theory)}, 12th issue (1992), Problem 12.86, and the 21st issue edition (2026), arXiv:1401.0300v48. The registry describes the collection as edited by E. I. Khukhro and V. D. Mazurov.
\item The registry labels this problem \emph{open} and records an update timestamp of \texttt{2026-10-16T14:00:00Z}. This is supplied registry metadata, not an independently verified literature assessment. The cited Notebook text has not been independently inspected for this draft.
\end{itemize}
\end{document}
